\subsection{Szczegóły implementacyjne algorytmów}

Symulacja działania algorytmów wymagała wiernego odwzorowania założeń teoretycznych z modelu Fornasiera, Schnass i Vybírala \cite{fornasier2012learning}. 

Zbiór punktów próbkowania $\mathcal{X} = \{\xi_1, \dots, \xi_{m_\mathcal{X}}\}$ na sferze jednostkowej $\mathbb{S}^{d-1}$ generowano poprzez losowanie wektorów z wielowymiarowego standardowego rozkładu normalnego $\mathcal{N}(0, I_d)$, a następnie normowanie każdego wektora do długości 1 (w normie $l_2$). Macierz pomiarowa $\Phi \in M_{m_\Phi \times d}$ została skonstruowana jako macierz losowa, której elementy przyjmowały wartości $+1$ lub $-1$ z równym prawdopodobieństwem (rozkład Bernoulliego), przeskalowane przez współczynnik $1/\sqrt{m_\Phi}$, aby zapewnić spełnienie Własności Ograniczonej Izometrii (RIP).

Pochodne kierunkowe funkcji $f$ estymowano poprzez ilorazy różnicowe z krokiem $\varepsilon$. Otrzymane wektory stanowiły prawą stronę układu równań $Y$, który następnie poddawano dekompozycji i optymalizacji.

\subsection{Weryfikacja dokładności rekonstrukcji (przypadek bezszumowy)}

W pierwszej fazie testów zbadano zachowanie algorytmów w wyidealizowanych warunkach dokładnych pomiarów. 